\documentclass[12pt,a4paper]{ctexart}

% ==================== 页面与基础设置 ====================
\usepackage[margin=1in]{geometry}
\usepackage{setspace}
\usepackage{indentfirst}
\usepackage{fancyhdr}
\usepackage{xcolor}

\setlength{\parindent}{2em}
\setlength{\headheight}{15pt}
\onehalfspacing

% ==================== 数学与定理 ====================
\usepackage{amsmath,amssymb,amsthm,bm}
\numberwithin{equation}{section}

\newtheorem{theorem}{定理}[section]
\newtheorem{lemma}[theorem]{引理}
\newtheorem{proposition}[theorem]{命题}
\newtheorem{definition}[theorem]{定义}
\newtheorem{remark}[theorem]{注}

% ==================== 图片、表格、算法 ====================
\usepackage{graphicx}
\usepackage{booktabs}
\usepackage{array}
\usepackage{multirow}
\usepackage{float}
\usepackage{caption}
\usepackage{enumitem}
\usepackage{mwe}
\usepackage{algorithm}
\usepackage{algpseudocode}

\captionsetup{font=small,labelfont=bf}

% ==================== 超链接 ====================
\usepackage{hyperref}
\hypersetup{
    colorlinks=true,
    linkcolor=blue,
    citecolor=blue,
    urlcolor=blue
}

% ==================== 页眉页脚 ====================
\pagestyle{fancy}
\fancyhf{}
\fancyhead[C]{后量子密码技术框架与典型方案分析}
\fancyfoot[C]{\thepage}

% ==================== 自定义命令 ====================
\newcommand{\keywords}[1]{\par\noindent\textbf{关键词：} #1}
\newcommand{\engkeywords}[1]{\par\noindent\textbf{Keywords:} #1}
\newcommand{\adv}{\mathrm{Adv}}

% ==================== 标题信息 ====================
\title{\heiti 后量子密码技术框架与典型方案分析}
\author{作者姓名\\[0.5em]
\small 某某大学~~某某学院}
\date{\today}

\begin{document}

\maketitle

% ==================== 中文摘要 ====================
\begin{abstract}
随着量子计算研究的快速发展，基于大整数分解与离散对数困难问题的经典公钥密码体制面临潜在威胁。后量子密码旨在构造即使在量子计算环境下仍具备可接受安全性的密码算法体系，已经成为密码学与网络安全领域的重要研究方向。本文以综述型技术论文的方式，对后量子密码的研究背景、基本安全模型、主要技术路线、典型标准化方案以及工程部署问题进行系统梳理。文中重点介绍格密码、码密码、多变量密码、哈希签名与同源密码等主要分支，并围绕 ML-KEM 与 ML-DSA 两类代表性方案分析其设计思想、核心流程和应用价值。在此基础上，本文进一步从安全依据、通信开销、密钥规模、实现复杂度和迁移可行性等角度对不同路线进行比较，讨论后量子密码从理论研究走向产业落地所面临的关键挑战。本文可作为后量子密码课程论文、基础综述与技术报告的写作框架。
\end{abstract}

\keywords{后量子密码；格密码；ML-KEM；ML-DSA；标准化；量子安全}

% ==================== 英文摘要 ====================
\vspace{1.5em}
\begin{center}
    {\bfseries\large Abstract}
\end{center}
With the rapid progress of quantum computing, classical public-key cryptosystems based on integer factorization and discrete logarithms face potential security threats. Post-quantum cryptography aims to construct cryptographic schemes that remain secure against adversaries equipped with quantum capabilities, and has become a major research direction in cryptography and cybersecurity. This paper presents a survey-style technical framework for post-quantum cryptography, covering its research motivation, security notions, major design families, representative standardized schemes, and deployment challenges. We focus on lattice-based, code-based, multivariate, hash-based, and isogeny-based approaches, with particular attention to ML-KEM and ML-DSA as representative standardized schemes. We further compare different approaches from the perspectives of security assumptions, communication overhead, key and signature sizes, implementation complexity, and deployment feasibility. The resulting manuscript serves as a practical foundation for course papers, introductory surveys, and technical reports on post-quantum cryptography.

\engkeywords{post-quantum cryptography, lattice-based cryptography, ML-KEM, ML-DSA, standardization, quantum security}

\tableofcontents
\newpage

% ==================== 正文 ====================
\section{引言}
公钥密码是现代网络安全体系的基础组成部分，广泛应用于密钥交换、数字签名、身份认证与安全通信等场景。当前工程上大量使用的 RSA、Diffie--Hellman 和椭圆曲线密码体制，主要依赖大整数分解问题或离散对数问题的计算困难性。然而，量子计算模型下的 Shor 算法表明，一旦大规模容错量子计算机实现，上述经典公钥体制的安全性将受到根本性冲击。

在这一背景下，后量子密码（Post-Quantum Cryptography, PQC）成为国际密码研究的热点。后量子密码并不要求密码系统运行在量子设备上，而是强调即使攻击者拥有量子计算能力，现有算法仍能在合理参数下维持安全性。与量子密钥分发等依赖特殊物理设备的方案不同，后量子密码更强调在现有数字基础设施上的软件与硬件可部署性，因此在协议升级与产业迁移方面具有现实可行性。

本文的目标并非给出某一具体算法的实现细节，而是构建一篇适合基础研究与课程论文写作的后量子密码综述型技术框架。本文主要贡献如下：
\begin{itemize}[leftmargin=2em]
    \item 梳理后量子密码的研究背景、基本概念与安全模型；
    \item 归纳后量子密码的主要技术路线及其代表性方案；
    \item 围绕 ML-KEM 与 ML-DSA 分析标准化方案的结构特点；
    \item 从安全、性能与部署角度比较不同后量子密码路线的优缺点。
\end{itemize}

\section{相关工作与标准化背景}
后量子密码的研究在过去二十余年中逐步形成了较为清晰的技术谱系。按困难问题来源划分，主流路线主要包括格密码、码密码、多变量密码、哈希签名以及同源密码等。不同路线分别基于不同的数学难题，因此在安全假设、密钥尺寸、计算开销和实现复杂度方面呈现出明显差异。

从国际标准化进程来看，美国国家标准与技术研究院（NIST）持续推动后量子密码标准化工作，使得研究重点逐步从纯理论安全性扩展到可实现性、互操作性和工程部署能力。当前，格密码路线因其在安全性、效率与通用性之间具备较好平衡，已经成为后量子公钥加密与数字签名标准化的核心方向。与此同时，Classic McEliece、SPHINCS+ 等方案也因其独特安全优势在长期研究中保持重要地位。

已有研究通常从三类角度展开。第一类关注数学基础与安全归约，如 LWE、RLWE 与 Module-LWE 问题的困难性分析；第二类关注具体密码方案的设计，例如基于模块格的密钥封装机制与签名体制；第三类关注工程实现与协议迁移，例如 TLS、VPN、PKI 等基础设施中引入后量子算法时面临的兼容与性能问题。本文在此基础上，更强调面向基础论文写作的整体技术框架与逻辑组织。

\section{基础概念与预备知识}
\subsection{后量子密码的基本定义}
\begin{definition}
后量子密码是指一类在经典计算环境中实现、但其安全性在面对具备量子计算能力的攻击者时仍被认为难以有效破坏的密码算法与协议体系。
\end{definition}

在功能上，后量子密码主要覆盖两类核心任务：一类是密钥建立与公钥加密，典型实现形式为密钥封装机制（Key Encapsulation Mechanism, KEM）；另一类是数字签名，用于消息完整性验证与身份认证。对于基础论文而言，围绕这两类功能组织全文结构，通常能够兼顾密码学理论与工程应用。

\subsection{安全模型与敌手优势}
为了描述后量子密码方案的安全性，通常需要先形式化定义敌手能力。对于密钥封装机制，常见目标是达到 IND-CPA 或 IND-CCA 安全；对于数字签名，常见目标是达到 EUF-CMA 安全。若记安全参数为 $\lambda$，敌手为 $\mathcal{A}$，则 KEM 的 IND-CCA 优势可记为
\begin{equation}
    \adv^{\mathrm{ind\mbox{-}cca}}_{\mathcal{A}}(\lambda)
    =
    \left|
    \Pr\bigl[\mathsf{Game}^{\mathrm{ind\mbox{-}cca}}_{\mathcal{A}}(\lambda)=1\bigr]
    - \frac{1}{2}
    \right|.
\end{equation}
若对任意多项式时间量子攻击者 $\mathcal{A}$，上述优势都可忽略，则称方案在该安全模型下是安全的。类似地，签名方案的不可伪造性也可写成相应的成功概率界，并通过归约证明与底层困难问题建立联系。

\subsection{格密码与 LWE 基本形式}
在当前后量子密码研究中，格密码路线占据核心地位。其原因在于格问题既具有较强的理论基础，又能支持加密、签名、同态计算等多类密码原语。设模数为 $q$，矩阵 $\bm{A}\in\mathbb{Z}_q^{m\times n}$，秘密向量为 $\bm{s}\in\mathbb{Z}_q^n$，误差向量为 $\bm{e}\in\mathbb{Z}_q^m$，则典型的模线性关系可写为
\begin{equation}
    \bm{t}=\bm{A}\bm{s}+\bm{e}\pmod q.
\end{equation}
该表达式刻画了学习带误差问题（Learning With Errors, LWE）的核心形式，即在存在小噪声干扰的情况下，从公开样本恢复秘密向量被认为是计算上困难的。进一步地，Module-LWE 可写为
\begin{equation}
    (\bm{A},\bm{B}=\bm{A}\bm{S}+\bm{E}) \in R_q^{k\times k}\times R_q^{k\times \ell},
\end{equation}
其中 $R_q$ 为模多项式环。基于这一结构，可以构造更高效的公钥加密、KEM 与签名方案。

\subsection{密码学环境示例}
\begin{definition}
设 $\Pi=(\mathsf{KeyGen},\mathsf{Encaps},\mathsf{Decaps})$ 为一个密钥封装机制。若算法 $\mathsf{KeyGen}$ 输出公钥与私钥对 $(pk,sk)$，$\mathsf{Encaps}(pk)$ 输出密文与会话密钥 $(c,K)$，且 $\mathsf{Decaps}(sk,c)$ 能够恢复正确的 $K$，则称 $\Pi$ 为一个正确的 KEM。
\end{definition}

\begin{theorem}
若底层 Module-LWE 问题在给定参数下对多项式时间量子攻击者是困难的，则基于该困难问题构造的标准化格密码 KEM 在适当变换下可达到 IND-CCA 安全。
\end{theorem}

\begin{proof}
该类结论通常通过一系列游戏跳转完成证明，将攻击者区分真实会话密钥与随机密钥的能力逐步归约到底层 Module-LWE 实例的可区分能力。对于基础论文而言，不必展开完整证明细节，但应明确安全性来自归约关系而非经验观察。
\end{proof}

\section{典型方案分析}
\subsection{后量子密码的主要技术路线}
从研究与部署角度看，后量子密码各路线的关注重点并不相同。格密码方案在效率与通用性上优势突出，适合构造 KEM 与签名；码密码历史悠久，安全基础较稳定，但公钥尺寸通常较大；多变量密码在签名方向曾受到关注，但部分候选方案经历了较强的密码分析；哈希签名依赖哈希函数安全性，结构简洁但状态管理或签名长度问题值得关注；同源密码具有优美数学结构，但在标准化竞赛中遭受较大挑战。

\subsection{ML-KEM 与密钥封装流程}
ML-KEM 是当前后量子密钥建立方向的代表性标准化方案，其设计继承了基于 Module-LWE 的格密码思想，并通过矩阵生成、噪声采样与压缩编码等机制实现较好的安全性与效率平衡。设方案参数为 $(q,k,\eta)$，则其公钥部分通常与公开矩阵和线性变换结果相关，而私钥部分包含恢复会话密钥所需的秘密信息。

从抽象层面看，ML-KEM 的正确性可表述为
\begin{equation}
    \Pr\bigl[\mathsf{Decaps}(sk,c)=K \mid (c,K)\leftarrow \mathsf{Encaps}(pk)\bigr] \ge 1-\varepsilon,
\end{equation}
其中 $\varepsilon$ 为解封装失败概率。为了保证协议层面的可靠性，设计者需要使该概率足够接近 $1$。同时，通信开销可粗略记为
\begin{equation}
    C_{\mathrm{kem}} = |pk| + |c| + |K|,
\end{equation}
该指标直接影响其在 TLS、VPN 与嵌入式设备中的部署成本。

\subsection{ML-DSA 与后量子签名}
ML-DSA 是后量子数字签名方向的重要代表方案，其设计同样建立在模块格问题之上，但重点转向短向量、挑战响应和拒绝采样等签名机制。与 KEM 相比，签名方案更关注签名大小、验证效率、随机性处理和侧信道防护。对工程系统而言，签名方案不仅影响通信负载，还直接影响证书体系、固件升级签名与代码签名等场景的可实施性。

若记签名大小为 $|\sigma|$，公钥大小为 $|pk|$，验证时间为 $T_{\mathrm{ver}}$，则可定义简单的部署代价函数
\begin{equation}
    \Gamma_{\mathrm{sig}} = \alpha |pk| + \beta |\sigma| + \gamma T_{\mathrm{ver}},
\end{equation}
其中 $\alpha,\beta,\gamma$ 为场景相关权重。该表达式虽然是抽象化指标，但有助于在综述型论文中统一讨论不同签名方案的实际使用代价。

\subsection{其他代表性方案}
除格密码方案外，Classic McEliece 在码密码方向具有较强长期安全声誉，但其公钥尺寸偏大；SPHINCS+ 作为无状态哈希签名方案，具有明确而稳健的安全基础，但签名与验证开销需要结合实际场景权衡。基础论文中通常不必对所有候选方案展开同等深度分析，但应明确这些路线在标准化与工程研究中的定位，以避免将后量子密码简单等同于格密码。

\section{后量子密码技术路线与典型流程}
\subsection{KEM 抽象流程}
从协议应用角度出发，后量子 KEM 通常可抽象为“生成密钥对、封装会话密钥、解封装恢复会话密钥”三步。设参与方为发送方与接收方，则典型流程如下：
\begin{algorithm}[H]
\caption{后量子 KEM 的基本流程}
\begin{algorithmic}[1]
\Require 安全参数 $\lambda$
\Ensure 会话密钥 $K$
\State 接收方执行 $(pk,sk)\leftarrow \mathsf{KeyGen}(\lambda)$
\State 接收方向发送方发布公钥 $pk$
\State 发送方计算 $(c,K)\leftarrow \mathsf{Encaps}(pk)$
\State 发送方向接收方发送密文 $c$
\State 接收方计算 $K' \leftarrow \mathsf{Decaps}(sk,c)$
\State 若 $K'=K$，则双方建立共享会话密钥
\end{algorithmic}
\end{algorithm}
该流程与经典密钥交换协议相比，在接口形式上较容易嵌入现有网络协议，因此成为后量子迁移中的关键组件。

\subsection{技术路线选择逻辑}
对于一篇基础技术论文而言，本节的作用不是给出新算法，而是说明为什么格密码会成为主线案例。其原因主要包括三点：一是底层问题具备较强理论研究基础；二是算法效率在标准化候选中表现较好；三是能够同时支持 KEM 与签名等多种功能。相比之下，其他路线往往在某一方面更突出，但整体平衡性略弱。

\section{理论分析与比较}
后量子密码方案的比较不能只停留在“是否量子安全”的层面，还需要综合考虑安全依据、通信开销、密钥规模、签名或密文长度、实现复杂度以及对现有协议栈的兼容性。为此，可定义一个抽象化的综合评价指标
\begin{equation}
    \Phi = \omega_1 S + \omega_2 E - \omega_3 C - \omega_4 M,
\end{equation}
其中 $S$ 表示安全性评价，$E$ 表示运行效率，$C$ 表示通信代价，$M$ 表示实现与维护复杂度，$\omega_i$ 为权重系数。虽然该指标不对应单一标准，但适合在综述型论文中统一组织分析维度。

若进一步考虑协议部署中的总通信成本，可将一次握手或认证的通信量写为
\begin{equation}
    B_{\mathrm{total}} = B_{\mathrm{base}} + B_{\mathrm{pq}} + B_{\mathrm{cert}},
\end{equation}
其中 $B_{\mathrm{pq}}$ 为引入后量子算法后新增的公钥、密文或签名字节开销。对资源受限设备而言，该项往往是实际部署时最敏感的因素之一。

从安全角度看，格密码的优势在于目前尚无与 Shor 算法同级别的量子攻击可以有效破坏大参数下的相关困难问题；但这并不意味着方案可直接无条件部署。参数选择、实现细节、采样偏差、解封装失败处理与侧信道防护都会影响实际安全性。因此，后量子密码的“安全”应被理解为“在当前已知数学与实现模型下的可信安全”，而非绝对安全。

\section{实验与结果展示}
\subsection{对比指标设计}
由于本文定位为基础综述型技术论文，不以具体代码实验为核心，因此实验部分采用表格和指标分析的方式展示不同路线的特征。设比较维度为安全基础、密钥规模、密文或签名大小、实现效率与适用场景，则可构造评价向量
\begin{equation}
    \bm{v}_i =
    \bigl[
    v_{i,1},
    v_{i,2},
    v_{i,3},
    v_{i,4},
    v_{i,5}
    \bigr]^{\top},
\end{equation}
其中 $\bm{v}_i$ 表示第 $i$ 类后量子方案的多维特征描述。进一步地，可通过加权方式得到总体评分
\begin{equation}
    Q_i = \sum_{j=1}^{5} w_j v_{i,j}, \qquad \sum_{j=1}^{5} w_j = 1.
\end{equation}
该类公式有助于为后续实验章节预留统一的分析接口。

\subsection{方案比较表}
表~\ref{tab:pqc-comparison} 给出了几类典型后量子密码路线的定性比较结果。
\begin{table}[H]
    \centering
    \caption{典型后量子密码路线对比}
    \label{tab:pqc-comparison}
    \begin{tabular}{p{2.8cm}p{2.5cm}p{2.2cm}p{2.2cm}p{2.4cm}}
        \toprule
        技术路线 & 安全基础 & 密钥/签名规模 & 计算效率 & 典型说明 \\
        \midrule
        格密码 & LWE / Module-LWE 等 & 中等 & 较高 & 当前标准化主流，兼顾 KEM 与签名 \\
        码密码 & 译码困难问题 & 公钥较大 & 较高 & 长期研究充分，Classic McEliece 代表性强 \\
        多变量密码 & 多变量方程求解困难 & 中等 & 依方案而异 & 曾受关注，但部分方案分析压力较大 \\
        哈希签名 & 哈希函数安全性 & 签名较大 & 中等 & 安全基础直接，SPHINCS+ 代表性较强 \\
        同源密码 & 椭圆曲线同源问题 & 较小到中等 & 较低 & 数学结构优美，但工程挑战较大 \\
        \bottomrule
    \end{tabular}
\end{table}

\subsection{结果分析}
从表~\ref{tab:pqc-comparison} 可以看出，格密码路线之所以在标准化中占据核心地位，主要是因为其在理论安全性、算法效率与功能多样性之间取得了较好平衡。特别是在 KEM 与签名两个方向上，格密码方案都表现出较强工程可行性。

然而，主流并不等于没有问题。首先，后量子算法的公钥、密文或签名通常比传统 ECC 方案更大，这会带来协议握手流量上升和证书体积增加的问题。其次，噪声采样、常数时间实现与故障注入防护等实现问题，会显著影响工程安全性。再次，后量子密码迁移并非孤立替换算法，而是需要与证书体系、硬件加速、协议协商机制和长期兼容策略协同推进。

\section{讨论}
从长期发展趋势看，后量子密码的研究已经从“寻找替代 RSA/ECC 的理论候选方案”转向“构建可部署、可验证、可迁移的完整密码生态”。这一转变意味着研究者不仅要关注底层困难问题和归约证明，还要关注协议适配、标准一致性、密钥管理、硬件实现与侧信道安全等系统层问题。

在工程迁移上，混合部署是一个现实可行的过渡策略。例如，在握手阶段同时结合经典密钥交换与后量子 KEM，有助于在保持兼容性的同时降低单一路线风险。与此同时，如何在性能敏感系统中控制通信膨胀、如何在受限设备中实现高质量随机采样，以及如何建立长期可信的参数更新机制，仍然是后量子密码落地中的关键议题。

\section{结论与未来工作}
本文围绕后量子密码构建了一套适用于基础研究与课程论文写作的综述型技术框架。文章首先说明了量子计算对经典公钥密码体制的潜在威胁，随后介绍了后量子密码的基本定义、安全模型与主要数学基础。在此基础上，本文重点分析了 ML-KEM 与 ML-DSA 等代表性标准化方案，并从安全性、效率、通信代价与工程迁移角度比较了不同技术路线的特点。

未来工作可以从三个方向进一步展开。第一，围绕具体方案补充更细致的安全证明、参数设计与实现细节；第二，结合实验平台给出更具体的性能测试与协议集成结果；第三，继续关注后量子密码标准化、混合部署与长期密钥基础设施升级等前沿问题。对于初学者而言，本文框架既可作为后量子密码基础论文的起点，也可作为后续深入研究单一算法路线的组织模板。

% ==================== 参考文献 ====================
\begin{thebibliography}{99}

\bibitem{nist-pqc}
NIST.
\newblock \textit{Post-Quantum Cryptography Project}.
\newblock Available online: \url{https://www.nist.gov/pqcrypto}.

\bibitem{fips203}
NIST.
\newblock \textit{Module-Lattice-Based Key-Encapsulation Mechanism Standard (ML-KEM)}.
\newblock FIPS 203, 2024.

\bibitem{fips204}
NIST.
\newblock \textit{Module-Lattice-Based Digital Signature Standard (ML-DSA)}.
\newblock FIPS 204, 2024.

\bibitem{regev}
Oded Regev.
\newblock On lattices, learning with errors, random linear codes, and cryptography.
\newblock \textit{Journal of the ACM}, 56(6), 2009.

\bibitem{kyber}
Joppe Bos, L\'eo Ducas, Eike Kiltz, Tancr\`ede Lepoint, Vadim Lyubashevsky, John M. Schanck, Peter Schwabe, Gregor Seiler, and Damien Stehl\'e.
\newblock CRYSTALS-Kyber: A CCA-secure module-lattice-based KEM.
\newblock In \textit{2018 IEEE European Symposium on Security and Privacy}, 2018.

\bibitem{dilithium}
L\'eo Ducas, Eike Kiltz, Tancr\`ede Lepoint, Vadim Lyubashevsky, Peter Schwabe, Gregor Seiler, and Damien Stehl\'e.
\newblock CRYSTALS-Dilithium: A lattice-based digital signature scheme.
\newblock \textit{IACR Transactions on Cryptographic Hardware and Embedded Systems}, 2018(1):238--268, 2018.

\bibitem{sphincs}
Daniel J. Bernstein, Andreas H\"ulsing, Stefan K\"olbl, Ruben Niederhagen, Joost Rijneveld, and Peter Schwabe.
\newblock SPHINCS+: Submission to the NIST Post-Quantum Project.
\newblock Available online.

\bibitem{mceliece}
Classic McEliece Team.
\newblock \textit{Classic McEliece: Conservative Code-Based Cryptography}.
\newblock Available online.

\end{thebibliography}

\end{document}
